\section{Output Derivation}\label{s:rollout}

Let $\boldsymbol r_z^{L_i}$ be the residual at the position of token $z \in \{d_1, d_2, s, o_1, o_2\}$ after layer $i$.

\subsubsection*{Layer 1}
The residuals at \texttt{SEP}, $o_1$ and $o_2$ after layer 1 are derived as follows:

\begin{equation}
\begin{split}
    \boldsymbol r_s^{L_1} &= \alpha_{s \to d_1}(\boldsymbol E_{d_1}+\boldsymbol P_{d_1}) + \alpha_{s \to d_2}(\boldsymbol E_{d_2}+\boldsymbol P_{d_2}) \\
    &\quad + \boldsymbol E_s + \boldsymbol P_s
\end{split}
\end{equation}
\begin{equation}
     \boldsymbol r_{o_i}^{L_1} = \boldsymbol E_m + \boldsymbol P_{o_i}, \quad i \in \{1,2\}.
\end{equation}
Note that we use our custom mask (outlined in Section \ref{s:method}) to prevent positions $4$ and $5$ from attending to anything in layer 1.

\subsubsection*{Layer 2}
The residuals at \texttt{SEP}, $o_1$ and $o_2$ after layer 2 are derived as follows:

\begin{equation}
    \boldsymbol r_{o_1}^{L_2} = \boldsymbol r_{o_1}^{L_1} + \beta_{o_1 \to s}\,\boldsymbol r_s^{L_1}
\end{equation}
\begin{equation}
    \boldsymbol r_{o_2}^{L_2} = \boldsymbol r_{o_2}^{L_1} + \beta_{o_2 \to s}\,\boldsymbol r_s^{L_1} + \beta_{o_2 \to o_1}\,\boldsymbol r_{o_1}^{L_1}
\end{equation}

\subsubsection*{Output logit at position 4 ($o_1$)}
We then derive the output logit at position 4 ($o_1$):

\begin{equation}
    \boldsymbol \ell_{o_1} = \boldsymbol r_{o_1}^{L_2}U = \big[\boldsymbol r_{o_1}^{L_1} + \beta_{o_1 \to s}\boldsymbol r_s^{L_1}\big]\boldsymbol U
\end{equation}

\begin{equation}
\label{eq:o1_logit_unsimp}
\begin{split}
    \boldsymbol \ell_{o_1} = 
    &\bigl[\boldsymbol E_m + \boldsymbol P_{o_1} 
    + \beta_{o_1 \to s} \bigl(
    \alpha_{s \to d_1}(\boldsymbol E_{d_1} + \boldsymbol P_{d_1}) \\
    &\qquad\quad
    + \alpha_{s \to d_2}(\boldsymbol E_{d_2} + \boldsymbol P_{d_2})
    + \boldsymbol E_s + \boldsymbol P_s
    \bigr)
    \bigr]\boldsymbol U
\end{split}
\end{equation}


\subsubsection*{Output logit at position 5 ($o_2$)}
Additionally, we derive the output logit at position 5 ($o_2$):

\begin{equation}
    \boldsymbol \ell_{o_2} = \boldsymbol r_{o_2}^{L_2}U = \big[\boldsymbol r_{o_2}^{L_1} + \beta_{o_2 \to s}\,\boldsymbol r_s^{L_1} + \beta_{o_2 \to o_1}\,\boldsymbol r_{o_1}^{L_1}\big]\boldsymbol U
\end{equation}

\begin{equation}
\label{eq:o2_logit_unsimp}
\begin{split}
    \boldsymbol \ell_{o_2} = 
    &\Bigl[\boldsymbol E_m + \boldsymbol P_{o_2} \\
    &\quad + \beta_{o_2 \to s} \Bigl(
        \alpha_{s \to d_1}(\boldsymbol E_{d_1} + \boldsymbol P_{d_1}) \\
        &\qquad + \alpha_{s \to d_2}(\boldsymbol E_{d_2} + \boldsymbol P_{d_2}) 
        + \boldsymbol E_s + \boldsymbol P_s
    \Bigr) \\
    &\quad + \beta_{o_2 \to o_1}(\boldsymbol E_m + \boldsymbol P_{o_1})
    \Bigr]\boldsymbol U
\end{split}
\end{equation}

\subsubsection*{Simplification}
We can simplify these formulae.
Let
\(
    \boldsymbol S = \boldsymbol E_s + \boldsymbol P_s
\).
Furthermore, since $o_1$ can only ever attend to $s$, we have $\beta_{o_1 \to s}=1$. Since $o_2$ can only attend to $s$ or $o_1$, we have $\beta_{o_2 \to s} + \beta_{o_2 \to o_1} = 1$. Therefore, we can simplify the equations \eqref{eq:o1_logit_unsimp} and \eqref{eq:o2_logit_unsimp} to get equations \eqref{eq:o1_logit_composition} and \eqref{eq:o2_logit_composition} in Section \ref{ss:results-circuits}.